\uuid{GHA9}
\exo7id{4765}
\titre{exo7 4765}
\auteur{quercia}
\organisation{exo7}
\datecreate{2010-03-16}
\isIndication{false}
\isCorrection{false}
\chapitre{Topologie}
\sousChapitre{Topologie des espaces vectoriels normés}
\cpge{mpsi-mp2i.espaces-vectoriels.a, pcsi.espaces-vectoriels.a, ptsi.geometrie-espace.c, mpsi-mp2i.applications-lineaires.a, pcsi.applications-lineaires.a, ptsi.applications-lineaires.a, mp-mpi.topologie-espaces-normes.a, pc.espaces-vectoriels-normes.a, psi.espaces-vectoriels-normes.a, mp-mpi.topologie-espaces-normes.c, pc.espaces-vectoriels-normes.c, psi.espaces-vectoriels-normes.c, mp-mpi.topologie-espaces-normes.e, pc.espaces-vectoriels-normes.e, psi.espaces-vectoriels-normes.e, pt.fonctions-deux-trois-variables.b}
\module{Analyse}
\niveau{L2}
\difficulte{}

\contenu{
\texte{
Soient $E,F$ deux evn et $G = E \times F$.
On pose pour $u = (\vec x,\vec y) \in G$ :
$$\|u\|_1 = \|\vec x\,\|_E + \|\vec y\,\|_F,\quad
  \|u\|_2 = \sqrt{\|\vec x\,\|_E^2 + \|\vec y\,\|_F^2},\quad
  \|u\|_\infty = \max(\|\vec x\,\|_E , \|\vec y\,\|_F).$$
}
\begin{enumerate}
    \item \question{Montrer que ce sont des normes sur $G$ et qu'elles sont deux {\`a} deux {\'e}quivalentes
    (sans hypoth{\`e}se de dimension finie).}
    \item \question{On prend $E=F$. Montrer que pour chacune de ces normes, l'application
    $G \to E, {(\vec x,\vec y\,)} \mapsto {\vec x+\vec y}$ est continue.}
\end{enumerate}
}
